\documentclass[preprint,12pt,authoryear]{elsarticle}
\usepackage[T1]{fontenc}
\usepackage[utf8]{inputenc}
\usepackage{amsmath,amssymb}
\usepackage{booktabs}
\usepackage{graphicx}
\usepackage{microtype}
\usepackage{hyperref}
\usepackage{float}
\hypersetup{hidelinks,hypertexnames=false}
\renewcommand{\thetable}{S\arabic{table}}
\renewcommand{\thefigure}{S\arabic{figure}}
\renewcommand{\theHtable}{S\arabic{table}}
\renewcommand{\theHfigure}{S\arabic{figure}}

\begin{document}
\begin{frontmatter}
\title{Supplementary Sensitivity Analyses for Positional-Parameter Robustness}
\author[aff1]{\DJ{}oko Ban{\dj}ur\corref{cor1}}
\ead{djoko.bandjur@pr.ac.rs}
\ead[url]{https://orcid.org/0000-0001-9034-6854}

\author[aff1]{Milo\v{s} Ban{\dj}ur}
\ead{milos.bandjur@pr.ac.rs}
\ead[url]{https://orcid.org/0009-0007-0124-3943}

\cortext[cor1]{Corresponding author}
\address[aff1]{Faculty of Technical Sciences, University of Priština, Kosovska Mitrovica, Serbia}
\end{frontmatter}

\section{Seed alignment and Welch sensitivity}

The ViT-B/16 and ViT-S/16 cohorts use the same six seed labels, but the
architectures consume different amounts of the global random stream during
model initialization. The training DataLoader does not use a separately
locked generator, so equal labels do not establish identical initial
weights or data order. Main-text intervals are seed-aligned difference
intervals retained for continuity with the within-architecture analyses.
Table~\ref{tab:s-welch} gives unpaired Welch intervals as a sensitivity
analysis.

\begin{table}[H]
\centering
\caption{Seed-aligned and unpaired Welch 95\% intervals for ViT-S minus
ViT-B nAUC contrasts. Both interval constructions give the same
interpretive conclusion for every family and metric.}
\label{tab:s-welch}
\small
\resizebox{\textwidth}{!}{%
\begin{tabular}{llrrr}
\toprule
PE family & Metric & Mean difference & Seed-aligned 95\% CI & Welch 95\% CI \\
\midrule
Learned & Adversarial nAUC & -0.218610 & $[-0.292924, -0.144295]$ & $[-0.293806, -0.143413]$ \\
Sinusoidal & Adversarial nAUC & -0.159924 & $[-0.184240, -0.135608]$ & $[-0.184361, -0.135487]$ \\
RoPE & Adversarial nAUC & -0.093204 & $[-0.160354, -0.026054]$ & $[-0.151529, -0.034878]$ \\
\midrule
Learned & Noise nAUC & -0.000894 & $[-0.003112, 0.001324]$ & $[-0.002435, 0.000646]$ \\
Sinusoidal & Noise nAUC & 0.000729 & $[-0.001757, 0.003216]$ & $[-0.001043, 0.002501]$ \\
RoPE & Noise nAUC & 0.000132 & $[-0.000479, 0.000743]$ & $[-0.000417, 0.000681]$ \\
\bottomrule
\end{tabular}%
}
\end{table}


\section{Ordinal robustness of the architecture-by-PE interaction}

Let $\Delta_f$ denote ViT-S minus ViT-B adversarial nAUC for family $f$.
More negative values indicate a larger loss under the smaller backbone.
Five of the six seeds have the order Learned $<$ Sinusoidal $<$ RoPE; seed
42 is the sole exception. Table~\ref{tab:s-orders} lists the six seed-level
orderings explicitly.

\begin{table}[H]
\centering
\caption{Per-seed ordering of adversarial architecture effects.}
\label{tab:s-orders}
\small
\resizebox{\textwidth}{!}{%
\begin{tabular}{rccc l}
\toprule
Seed & Learned $\Delta$ & Sinusoidal $\Delta$ & RoPE $\Delta$ & Ascending order \\
\midrule
42 & -0.121191 & -0.159980 & -0.153736 & Sinusoidal $<$ RoPE $<$ Learned \\
123 & -0.188993 & -0.131088 & -0.065368 & Learned $<$ Sinusoidal $<$ RoPE \\
456 & -0.200771 & -0.197074 & -0.157989 & Learned $<$ Sinusoidal $<$ RoPE \\
789 & -0.205311 & -0.174361 & 0.014230 & Learned $<$ Sinusoidal $<$ RoPE \\
1011 & -0.326862 & -0.149978 & -0.087133 & Learned $<$ Sinusoidal $<$ RoPE \\
1213 & -0.268529 & -0.147064 & -0.109227 & Learned $<$ Sinusoidal $<$ RoPE \\
\bottomrule
\end{tabular}%
}
\end{table}


The rank sums are 8, 11, and 17 for Learned, Sinusoidal, and RoPE,
respectively, reproducing the exact Friedman result $Q=7.0$,
$p=0.0289352$. Table~\ref{tab:s-did} reports the corresponding
pairwise difference-in-differences contrasts.

\begin{table}[H]
\centering
\caption{Pairwise differences in architecture effects. Same-sign counts
refer to the dominant sign across the six seeds.}
\label{tab:s-did}
\small
\resizebox{\textwidth}{!}{%
\begin{tabular}{lrr}
\toprule
Difference-in-differences contrast & 95\% CI & Same-sign seeds \\
\midrule
Learned $-$ RoPE & $[-0.235330, -0.015482]$ & 5/6 negative \\
Sinusoidal $-$ RoPE & $[-0.133297, -0.000143]$ & 6/6 negative \\
Learned $-$ Sinusoidal & $[-0.141594, 0.024222]$ & 5/6 negative \\
\bottomrule
\end{tabular}%
}
\end{table}


\section{Architecture-dependent seed dispersion}

Table~\ref{tab:s-dispersion} reports the seed-level dispersion of the
AMP-matched architecture effects.

\begin{table}[H]
\centering
\caption{Seed dispersion of adversarial nAUC. The leave-one-out range is
the minimum and maximum ViT-S/ViT-B SD ratio after omitting each seed in turn.}
\label{tab:s-dispersion}
\small
\resizebox{\textwidth}{!}{%
\begin{tabular}{lrrrrrr}
\toprule
PE family & ViT-B range & ViT-S range & ViT-B SD & ViT-S SD & SD ratio & LOO SD-ratio range \\
\midrule
Learned & 0.007196 & 0.209900 & 0.002436 & 0.071663 & 29.414 & $[19.930, 47.272]$ \\
RoPE & 0.035298 & 0.149260 & 0.011945 & 0.055722 & 4.665 & $[3.919, 5.754]$ \\
Sinusoidal & 0.051913 & 0.059656 & 0.018077 & 0.019832 & 1.097 & $[0.614, 1.485]$ \\
\bottomrule
\end{tabular}%
}
\end{table}


The dispersion ordering, Learned $\gg$ RoPE $>$ Sinusoidal, differs from
the mean architecture-effect ordering, Learned $>$ Sinusoidal $>$ RoPE.
This is reported descriptively; no mechanism or universal variance-scaling
law is inferred.

\section{Noise sign balance and normalized values above one}

Table~\ref{tab:s-signs} gives the seed-level sign balance for the noise
architecture contrasts.

\begin{table}[H]
\centering
\caption{Seed signs for ViT-S minus ViT-B nAUC contrasts.}
\label{tab:s-signs}
\begin{tabular}{lrrrr}
\toprule
PE family & Attack $\Delta<0$ & Attack $\Delta>0$ & Noise $\Delta<0$ & Noise $\Delta>0$ \\
\midrule
Learned & 6 & 0 & 3 & 3 \\
Sinusoidal & 6 & 0 & 2 & 4 \\
RoPE & 5 & 1 & 3 & 3 \\
\bottomrule
\end{tabular}
\end{table}


Noise contrasts have no common sign tendency, unlike the adversarial
contrasts. ViT-S Sinusoidal seed 123 has noise nAUC 1.0009374. Normalized
accuracy is defined relative to clean accuracy on the same finite split
and is not clipped at one; a perturbed model can therefore marginally
outperform its clean reference, and its integrated normalized curve can
slightly exceed one.


\section{Sensitivity to the adversarial lower envelope}
\label{sec:s-no-envelope}

The primary analysis retains the highest-loss restart at each native budget
and then applies a cumulative lower envelope to adversarial normalized
accuracy after ordering points by achieved $\rho_{\mathrm{rel}}$. As a
deterministic sensitivity analysis, we keep the same selected restart at
every budget but connect the raw measured normalized accuracies directly,
using the same exact clean anchor, achieved-displacement ordering,
piecewise-linear interpolation, and no clipping above one.

\begin{table}[H]
\centering
\caption{Deterministic sensitivity to omitting the adversarial cumulative
lower envelope. Values are family-mean noise-minus-adversarial nAUC
gaps across six seeds. ``Raw'' connects the same loss-selected
native-budget points without enforcing monotonicity; $\Delta$ is raw
minus envelope. Every raw seed-level gap is positive at both endpoints.}
\label{tab:s-no-envelope}
\scriptsize
\resizebox{\textwidth}{!}{%
\begin{tabular}{llrrr rr r c}
\toprule
 & & \multicolumn{3}{c}{$\rho_{\max}=0.09$} & \multicolumn{3}{c}{$\rho_{\max}=0.095$} & \\
\cmidrule(lr){3-5}\cmidrule(lr){6-8}
Dataset & PE family & Envelope gap & Raw gap & $\Delta$ & Envelope gap & Raw gap & $\Delta$ & Raw gap $>0$ \\
\midrule
ImageNet-100 & Learned & 0.004501 & 0.004466 & -0.000036 & 0.004802 & 0.004768 & -0.000034 & 6/6 \\
 & Sinusoidal & 0.137345 & 0.137345 & +0.000000 & 0.166796 & 0.166794 & -0.000001 & 6/6 \\
 & RoPE & 0.047449 & 0.047397 & -0.000053 & 0.056335 & 0.056285 & -0.000050 & 6/6 \\
 & ALiBi & 0.126587 & 0.126587 & +0.000000 & 0.140466 & 0.140466 & +0.000000 & 6/6 \\
\midrule
CIFAR-100 & Learned & 0.015369 & 0.014591 & -0.000778 & 0.018347 & 0.017030 & -0.001317 & 6/6 \\
 & Sinusoidal & 0.040536 & 0.040536 & +0.000000 & 0.044829 & 0.044829 & +0.000000 & 6/6 \\
 & RoPE & 0.008677 & 0.008540 & -0.000138 & 0.010009 & 0.009879 & -0.000130 & 6/6 \\
 & ALiBi & 0.020276 & 0.020276 & +0.000000 & 0.022697 & 0.022697 & +0.000000 & 6/6 \\
\bottomrule
\end{tabular}%
}
\end{table}


At both integration endpoints, all 48 raw noise-minus-adversarial
seed-level gaps remain positive and every within-seed family rank
configuration is preserved. Consequently, the exact Friedman results are
unchanged: $Q=16.4$, exact $p=8.03755\times10^{-6}$ for ImageNet-100 and
$Q=15.0$, exact $p=9.46924\times10^{-5}$ for CIFAR-100. The largest change
at the primary endpoint is CIFAR-100 Learned PE, whose family-mean gap
changes from $0.015369$ with the envelope to $0.014591$ without it. At
$\rho_{\max}=0.095$, the largest change is again CIFAR-100 Learned PE,
from $0.018347$ to $0.017030$. Thus, the positive-gap and between-family
heterogeneity conclusions do not depend on the lower-envelope operation.
Because the no-envelope analysis connects the same selected points without a
cumulative minimum, it also removes the family-dependent opportunity for
envelope selection that can arise when curves contain unequal numbers of
measured adversarial budgets.


\section{Sensitivity to the post-lock task-loss budget}
\label{sec:s-postlock-grid}

The canonical ViT-S/16 Sinusoidal task-loss curves were defined on the
prespecified native-budget grid ending at $0.010$. A historical targeted
coverage extension subsequently attempted to reach $\rho_{\mathrm{rel}}=0.09$
for seeds 123, 456, and 1213 using additional native budgets through 0.020; the
run ledger records all three as target-not-reached and the registered combined
maxima equal their canonical maxima to ledger precision. Original row-level
coverage-extension JSONs are not redistributed. The later all-restart audit is
a separate public row-level audit that used the same checkpoints, split,
float32 tensors/model state without autocast or gradient scaling, task-loss
objective, and $200\times5$ schedule, and additionally evaluated native budget
$0.020$ for all six seeds. Because that point was collected after the canonical
grid was locked, it is not inserted into the primary estimand. The analogous
ViT-B/16 budget-0.015 audit points reach
$\rho_{\mathrm{rel}}\in[0.096597,0.106636]$, above both the locked
ViT-B/16 endpoint and the $0.095$ sensitivity endpoint, so they cannot alter
the reported ViT-B/16 nAUC values.

As a transparent post-lock sensitivity, we add the loss-selected $0.020$
point to each of the six ViT-S/16 Sinusoidal curves and retain the same
cross-architecture integration endpoint
$\rho_{\mathrm{rel}}=0.0792384$.

\begin{table}[H]
\centering
\caption{Post-lock sensitivity to adding the loss-selected ViT-S/16
Sinusoidal task-loss point at native budget 0.020. The primary column
uses the prespecified canonical grid ending at 0.010; the sensitivity
column adds the later all-restart audit point while keeping the locked
cross-architecture endpoint $\rho_{\mathrm{rel}}=0.0792384$.}
\label{tab:s-postlock-grid}
\scriptsize
\resizebox{\textwidth}{!}{%
\begin{tabular}{rrrrrr}
\toprule
Seed & Audit $\rho_{\mathrm{rel}}$ & Audit normalized accuracy & Primary adversarial nAUC & With 0.020 point & Difference \\
\midrule
42 & 0.079453 & 0.021556 & 0.796951 & 0.790104 & -0.006847 \\
123 & 0.071233 & 0.026967 & 0.779643 & 0.732882 & -0.046760 \\
456 & 0.068338 & 0.018950 & 0.725753 & 0.687994 & -0.037759 \\
789 & 0.066990 & 0.025229 & 0.812803 & 0.702306 & -0.110496 \\
1011 & 0.076631 & 0.020154 & 0.766697 & 0.752726 & -0.013971 \\
1213 & 0.067808 & 0.026297 & 0.716350 & 0.664651 & -0.051699 \\
\midrule
\multicolumn{3}{r}{Across-seed mean $\pm$ sample SD} & $0.766366\pm0.038521$ & $0.721777\pm0.045899$ & $-0.044589$ \\
\bottomrule
\end{tabular}%
}
\end{table}


The added point lowers adversarial nAUC in every seed. The across-seed mean
changes from $0.766366\pm0.038521$ to
$0.721777\pm0.045899$, a difference of $-0.044589$. This increases the
noise-minus-adversarial separation and leaves the adversarial family
ordering unchanged. The sensitivity therefore strengthens the central
decoupling result, while the prelocked-grid values remain the primary
reported estimates.

\end{document}
